\section{Design philosophy and invariants}\label{sec:philosophy}

The part tree is the most widely shared piece of state in QtRocket: the ODE loop reads its
mass properties per integrator stage (\cref{sec:flight}), the serializer walks it
(\cref{sec:serialization}), the CLI addresses into it (\cref{sec:cli}), and the Qt tree view
mirrors it row-for-row (\cref{sec:gui}). A structure with that many consumers cannot rely on
discipline; it needs invariants that hold \emph{by construction}. This section states the
eight commitments that shape every type and function in the design. Each one is chosen to move
a class of defect from ``a bug someone must remember not to write'' to ``a state the program
cannot represent'' --- or, where full unrepresentability is impossible, to a hard, located
failure instead of a silently wrong number.

The commitments are stated here as rules, before any type is introduced, because their value
is precisely that they constrain everything downstream: the leaf contract
(\cref{sec:leaf}), the tree and its verbs (\cref{sec:tree}), the cache gating
(\cref{sec:caching}), the motor slot (\cref{sec:motor}), and the GUI adapter
(\cref{sec:gui}) are all forced to be mutually consistent with them.

\begin{designrule}[The eight invariants]
\textbf{1. One mutation authority.} Reads hand out \code{const Part\&}; every post-attach
write is a \code{PartsModel} verb that writes and then invalidates the correct cache chain. \\[2pt]
\textbf{2. Linear ownership.} Nodes own their \code{Part} and their children by
\cppinline|std::unique_ptr|; double-attach, cycles, and retained aliases are uncompilable. \\[2pt]
\textbf{3. Derived state is gated.} Caches rebuild iff their inputs moved, and a gate is
stamped only after a successful build. \\[2pt]
\textbf{4. One id space.} \code{Part::Id} --- process-unique, fresh on clone, never
serialized --- is the only identity that survives across edits. \\[2pt]
\textbf{5. Typed errors.} Structural verbs return \cppinline|std::expected|; a failed
mutation is a value the caller can switch on, never a silent no-op. \\[2pt]
\textbf{6. ``No design'' is representable.} The root is nullable and \code{hasDesign()}
reports it; there is no placeholder part. \\[2pt]
\textbf{7. Single-writer by contract.} The \code{const} accessors memoize through mutable
caches; all access comes from one thread at a time, and the contract is documented where it
can be broken. \\[2pt]
\textbf{8. Events are aboutTo/did pairs.} Every mutation fires one notification before and
one after the tree changes, carrying the row coordinates Qt's
\code{beginInsertRows}/\code{endInsertRows} contract requires.
\end{designrule}

\subsection{One authority, one owner}\label{sec:philosophy-authority}

The first two invariants are one idea seen from two sides: \emph{who may write} and
\emph{who may hold}. \code{PartsModel} is the single mutation authority
(\src{core/model/PartsModel.h}{133}): reads hand out \code{const Part\&}, and the only
post-attach write paths are its routed verbs --- \code{setPartMass}, \code{setPartInertia},
\code{setLink}, \code{mutatePart}, the structural \code{attach}/\code{detach} family, and the
motor verbs (\cref{sec:tree}). Each verb performs the plain write and then dirties exactly
the cache chains that write can invalidate. The leaf cooperates by construction: \code{Part}'s
setters are \emph{private}, exposed only through a single friend seam.

\begin{listing}[htbp]
\begin{minted}{cpp}
private:
    friend class model::PartsModel;  // the sole post-attach mutation seam

    /// Plain writes: PartsModel's routed verbs own cache invalidation.
    void setMass(double m) { mass = m; }
    void setI(const Matrix3& I) { inertiaTensor = I; }
\end{minted}
\caption{The mutation seam on the leaf (\srccap{core/model/parts/Part.h}{133}). The setters
carry no invalidation duty and no virtual dispatch; routing and invalidation are entirely the
authority's job.}\label{lst:philosophy-seam}
\end{listing}

\begin{keyinsight}[Invalidation cannot be skipped, by construction]
Cache invalidation is the classic ``must remember'' duty: one forgotten dirty-flag write and a
composite tensor is silently stale for the rest of the session. Here the compiler carries the
duty. Client code cannot reach \code{setMass} at all --- the only path to a leaf write runs
through a \code{PartsModel} verb, and every verb pairs the write with its invalidation
(\src{core/model/PartsModel.cpp}{399}). The one deliberately open door, \code{mutatePart},
requires the caller to declare a \code{MutateHint} (\src{core/model/PartsModel.h}{139}) so
that even the escape hatch states which chains its edit can dirty --- a mass-only edit never
re-runs the geometry sweep.
\end{keyinsight}

Ownership is equally linear. Each \code{PartNode} owns its \code{Part} and its children
through \cppinline|std::unique_ptr| (\src{core/model/PartsModel.h}{108}); \code{PartsModel}
owns the root; \code{RocketModel} owns the model. A part therefore has exactly one owner at
every instant of its life, and the transfer points are visible in the signatures:
\code{attach} \emph{consumes} a \cppinline|std::unique_ptr<part::Part>|, and \code{detach}
\emph{returns} the owned sub-tree to the caller as a first-class value
(\src{core/model/PartsModel.h}{216}) --- the natural undo and clipboard memento. Duplication
has exactly one spelling, \code{clone()}, and the assignment operators are deleted so a
subclass can never be sliced into a base \code{Part} (\src{core/model/parts/Part.h}{56}).

\begin{rejected}[Rejected --- shared ownership with runtime membership checks]
The alternative is a \cppinline|std::shared_ptr| forest in which the tree checks, at attach
time, that a part is not already attached elsewhere. This demotes every ownership violation
from a compile error to a runtime bug class: double-attach becomes a check that must be
written (and reached), a cycle becomes a traversal that must terminate, and --- worst --- a
\emph{retained alias} remains perfectly legal: any client that keeps its
\cppinline|shared_ptr| after attaching can write through it forever, past every routed verb,
producing precisely the unsanctioned mutation invariant~1 exists to eliminate. Under
\cppinline|unique_ptr| the retained alias is uncompilable rather than unsanctioned. The
mass-delta gate of \cref{sec:caching} remains as a backstop for any residual
\cppinline|const_cast|-class bypass, but the front door is closed by the type system, not by
review.
\end{rejected}

\subsection{Derived state is gated}\label{sec:philosophy-gates}

Composite mass properties, resolved placements, and the overlap verdict are all
\emph{derived} state: pure functions of the leaves, the links, and (for mass) the query time.
Invariant~3 says derived state may be cached only behind a gate that provably tracks its
inputs. The design has exactly two gates, one per kind of input motion: the \emph{mass-delta
gate} rebuilds the composite CM and tensor when the live mass sum moves (a burning motor) or
a mass edit lands, and the \emph{placement gate} re-resolves geometry and re-runs the envelope
sweep only when a node or link changes --- geometry is invariant under a burn, so the hot ODE
path never pays for a re-resolve. \Cref{sec:caching} develops both gates in full.

What belongs to the philosophy is the stamping rule, because it is a correctness property and
not an optimization:

\begin{keyinsight}[Stamps land only after success]
\code{ensureCompositeCache} computes the rebuild \emph{first} and writes the gate stamps
(\cppinline|builtAtMass_|, clearing \cppinline|massDirty_|) only after
\code{computeCompositeAt} returns (\src{core/model/PartsModel.cpp}{122}). The ordering is
load-bearing: a self-intersecting design makes the composite build throw
(\src{core/model/PartsModel.cpp}{184}), and because nothing was stamped, the throw recurs on
\emph{every} read. A cache that recorded the failed build as done would convert a hard,
located failure into a stale tensor served indefinitely --- the worst outcome available to a
simulator, a plausible wrong number.
\end{keyinsight}

\subsection{One id space}\label{sec:philosophy-id}

Every consumer that must refer to a part across edits --- GUI selection, CLI commands, event
payloads, the model's own index --- uses the same handle: \code{Part::Id}, a
\cppinline|std::uint64_t| minted from a process-wide atomic counter at construction
(\src{core/model/parts/Part.cpp}{16}). The counter starts at 1 so that 0 is reserved as
``none/invalid,'' which the event payloads exploit for roots and clearing resets
(\src{core/model/PartsModel.h}{147}). Three properties make one id space sufficient:

\begin{itemize}
  \item \textbf{Fresh on clone.} Every constructor mints a new id, including the protected
  copy constructor that backs \code{clone()} (\src{core/model/parts/Part.cpp}{37}). Two parts
  never share an id, however identical their names and mass properties.
  \item \textbf{Never serialized.} Ids are meaningful only within a single run
  (\src{core/model/parts/Part.h}{113}); a loaded design mints fresh ids. Persisting them
  would import cross-file collision handling into every load path for no benefit --- nothing
  outside the process ever needs the handle.
  \item \textbf{Fails safe when stale.} \code{PartsModel::find} is an $O(1)$ id-keyed lookup
  that returns \code{nullptr} for an absent id (\src{core/model/PartsModel.h}{167}). A
  consumer holding the id of a detached part gets a testable miss, where a retained pointer
  would dangle.
\end{itemize}

\begin{designrule}[Names are labels; ids are identity]
Part names are human-facing and deliberately non-unique; \code{findByName} exists for
console convenience and is documented as pre-order first-match
(\src{core/model/PartsModel.h}{179}). No correctness property anywhere in the system hangs
on a name. There is also no second handle type --- no node pointer or row index offered as an
identity --- because every additional handle kind is one more thing a consumer can hold
stale. One id space means one staleness story, and it is fail-safe.
\end{designrule}

\subsection{Typed errors, and a mutation that cannot silently fail}\label{sec:philosophy-errors}

Structural verbs return \cppinline|std::expected|: \code{attach} and \code{attachSubtree}
yield the new id or an \code{AttachError} (\code{NullPart}, \code{NoSuchParent},
\code{DuplicateMotor}); \code{detach} yields the owned sub-tree or a \code{DetachError}
(\code{NoSuchId}, \code{IsRoot}) (\src{core/model/PartsModel.h}{136}). The leaf-edit verbs,
whose only failure mode is an unknown id, return \code{bool}. In every case failure is a
value at the call site: the CLI switches on the error code and prints the precise cause
(\src{cli/Repl.cpp}{957}), and a test can assert the exact rejection rather than probing for
side effects.

\begin{rejected}[Rejected --- logged void no-ops]
The tempting minimal interface --- \code{void attach(...)} that logs and returns on failure
--- makes failure invisible at the call site. A caller that must know whether the attach
happened is reduced to diffing the tree's size before and after, and even then the
\emph{cause} is gone: was the parent id wrong, the part null, the motor slot occupied? The
log knows; the program does not. Errors-as-values costs one \cppinline|std::expected| in the
signature and buys call sites that are impossible to write incorrectly by accident: for
\code{detach}, the success value \emph{is} the detached sub-tree, so a caller cannot take the
value while ignoring the possibility of failure --- they arrive in the same object.
\end{rejected}

The \code{DuplicateMotor} rejection illustrates why the errors are typed rather than boolean:
the single-motor invariant is enforced at the doors through which a part can join an owned
tree --- \code{attach} and \code{attachSubtree}, which scans the entire incoming sub-tree
(\src{core/model/PartsModel.cpp}{319}) --- and the caller is told
\emph{that} the motor slot is the objection, not merely that ``something'' failed
(\cref{sec:motor}).

\subsection{The empty state is real}\label{sec:philosophy-empty}

A freshly started session, or one whose design was cleared, has no rocket. The design
represents that state honestly: the root is nullable
(\src{core/model/PartsModel.h}{259}), \code{hasDesign()} reports it
(\src{core/model/PartsModel.h}{161}), \code{size()} is 0, and every query fails safe ---
\code{find} misses, the pre-order walk visits nothing. Consumers respond in their own idiom;
the CLI's contract is an explicit \code{ERR \ldots: no design} on commands that need one
(\cref{sec:cli}).

\begin{rejected}[Rejected --- a mandatory placeholder part]
The alternative is to guarantee a non-null root by installing a placeholder part in an empty
model. This trades an honest null for a dishonest object: the placeholder has \emph{some}
mass, \emph{some} geometry, and a simulation will happily consume both --- flying a rocket
that does not exist and returning plausible numbers for it. Every consumer must then ask
``is this the placeholder?'' before trusting the tree, which is exactly the null check again,
now spelled as a magic-value comparison with strictly worse failure modes when someone
forgets it. \code{hasDesign()} is the same predicate with a compiler-visible spelling and no
physics attached.
\end{rejected}

\subsection{Single-writer by contract}\label{sec:philosophy-threading}

The composite accessors are \code{const} but not concurrently callable, and the design says
so out loud rather than letting the signatures imply otherwise. \code{compositeCm},
\code{compositeI}, \code{placementDiagnostics}, and \code{resolvedPlacements} all memoize
through \code{mutable} caches (\src{core/model/PartsModel.h}{114}), and the motor's log-once
burnout latch writes inside \code{const} \code{thrust(t)}
(\src{core/model/MotorModel.h}{342}). Concurrent access --- even all-\code{const} --- is
therefore a data race outright. The contract is single-writer: all access from one thread at
a time, stated in the class documentation and again at the cache fields
(\src{core/model/PartsModel.h}{39}), so that any future threading work reopens the question
deliberately instead of trusting \code{const}.

\begin{keyinsight}[The race is deeper than torn reads]
Even a hypothetically race-free interleaving would be wrong: the composite CM and the
composite tensor are only meaningful as a \emph{pair} from the same rebuild --- the tensor is
expressed about that CM. Two threads interleaving through the gate could pair a CM from one
rebuild with a tensor from another: two individually valid values forming a physically
meaningless answer. Any concurrency story must therefore preserve multi-value consistency,
not merely atomicity of individual reads.
\end{keyinsight}

The recorded direction, should a real concurrency requirement appear, is
\emph{snapshotting}: each thread flies its own \code{PartNode::clone()} and only the handoff
synchronizes. \Cref{sec:tradeoffs-threading} develops that plan and the rejected locking
alternative.

\subsection{Events shaped for the strictest subscriber}\label{sec:philosophy-events}

Every mutation --- structural, link, leaf edit, motor swap, and whole-tree reset --- fires a
change notification as an \emph{aboutTo/did pair}: the same \code{Event} value delivered
once with \code{before == true} immediately before the tree changes and once with
\code{before == false} immediately after (\src{core/model/PartsModel.h}{143}). The payload
is \code{\{kind, id, parentId, row\}} --- the affected node, its (ex-)parent, and its row
within that parent.

That shape is dictated by the most demanding subscriber the model must serve. Qt's
\code{QAbstractItemModel} contract requires \code{beginInsertRows}/\code{beginRemoveRows}
to be called \emph{before} the underlying data changes, with parent-and-row coordinates, and
the matching \code{end} call after. Because the event stream fires exactly that way, the GUI
adapter is a one-to-one mapping with no bookkeeping of its own: \code{aboutTo} maps to
\code{begin*}, \code{did} maps to \code{end*}, and the payload supplies the coordinates
verbatim (\src{gui/RocketTreeView.cpp}{29}; \cref{sec:gui}). Subscribers with weaker needs
--- a view that repaints wholesale --- simply ignore the \code{before} leg. The registration
is a single callback slot rather than a multicast list
(\src{core/model/PartsModel.h}{248}): the model layer stays Qt-free, one subscriber exists,
and the callback fires synchronously on the mutating thread, which the single-writer
contract already pins.

\medskip

None of these invariants is free-standing: the leaf's private seam (\cref{sec:leaf}) exists
to serve invariant~1, the verbs of \cref{sec:tree} are shaped by invariants~1, 2, and~5, the
gates of \cref{sec:caching} are invariant~3 made concrete, and the trade-off discussion of
\cref{sec:tradeoffs} weighs what each rule costs. The test suite exercises them directly ---
rejected attaches, stale-id misses, the empty state, and the after-success stamping rule are
all asserted behaviors, not aspirations (\cref{sec:verification}).
